\section{支撑材料清单}

本文全部数值结果由下列 Python 源程序真实计算生成，入口为 \verb|main.py|，
运行环境为 Python 3 与 NumPy、SciPy、pandas、openpyxl。程序不含任何硬编码的结果数值，
四个子问题的答案均由同一求解核在不同物性闭合与终止判据下计算得到。

\begin{longtable}{lp{5.0cm}p{6.2cm}}
  \caption{支撑材料文件清单}\label{tab:files} \\
  \toprule
  类别 & 文件名 & 内容 \\
  \midrule
  \endfirsthead
  \toprule
  类别 & 文件名 & 内容 \\
  \midrule
  \endhead
  \bottomrule
  \endlastfoot
  程序入口 & \verb|main.py|          & 四问求解编排与结果写出 \\
           & \verb|driver.py|        & 时间推进、算子分裂与 Picard 迭代 \\
           & \verb|fvm.py|           & 柱坐标有限体积离散与三对角求解 \\
           & \verb|props.py|         & 附录2--4 物性闭合式 \\
           & \verb|ambient.py|       & 附件1 环境温湿度插值 \\
           & \verb|geometry.py|      & 附件2 半径单调化与 PCHIP 插值 \\
           & \verb|params.py|        & 题给参数与数据读入 \\
           & \verb|problem1.py|      & 问题一预热段求解与报表 \\
           & \verb|problem2.py|      & 问题二全程变物性求解与报表 \\
           & \verb|problem3.py|      & 问题三达标时长求解与报表 \\
           & \verb|problem4.py|      & 问题四收缩域求解、效应隔离与干物质核算 \\
           & \verb|master_p23.py|    & 问题二、三共用的主求解流程 \\
           & \verb|validate.py|      & 守恒、收敛、级数解对照与极限退化检验 \\
           & \verb|io_out.py|        & 报表与工作簿写出 \\
           & \verb|data_check.py|    & 附件数据完整性核对 \\
           & \verb|template_check.py|& 附件3 模板结构核对 \\
  \midrule
  计算结果 & \verb|result1.xlsx|     & 问题一温度与水分浓度全网格结果 \\
           & \verb|result2.xlsx|     & 问题二全烘干过程双表结果 \\
           & \verb|result3.xlsx|     & 问题三 $0\to t_{\rm end}$ 结果 \\
           & \verb|result4.xlsx|     & 问题四收缩域结果 \\
           & \path{final_diagnostics/q4_effect_isolation_final.csv} & 情景 S0--S4 达标时长 \\
           & \verb|mass_consistency.csv| & 干物质守恒诊断 \\
           & \path{final_diagnostics/q3_sensitivity_final.csv} & 单因子灵敏度记录 \\
           & \verb|validation.json|  & 数值可信度检验记录 \\
\end{longtable}

\section{程序源代码}

下列代码为完整可运行源程序，按依赖层次排列：先给出参数与数据读入，再给出离散核与物性、
时间推进与求解编排，最后给出四问的调用与验证。运行方式为在 \verb|code/| 目录下执行
\verb|python main.py|，程序自动读入附件并把全部报表与工作簿写入 \verb|output/|。

\subsection{参数与数据读入：params.py}
\lstinputlisting{../code/params.py}

\subsection{环境驱动插值：ambient.py}
\lstinputlisting{../code/ambient.py}

\subsection{收缩半径插值：geometry.py}
\lstinputlisting{../code/geometry.py}

\subsection{物性闭合式：props.py}
\lstinputlisting{../code/props.py}

\subsection{有限体积离散核：fvm.py}
正式问题三采用 $N=2560$、$\Delta t=0.25$~s，首达时间为 $57.4716$~h；
界面采用正文定义的算术平均口径。
\lstinputlisting[linerange={1-32,42-10000}]{../code/fvm.py}

\subsection{时间推进与耦合迭代：driver.py}
\lstinputlisting{../code/driver.py}

\subsection{问题一求解：problem1.py}
\lstinputlisting{../code/problem1.py}

\subsection{问题二、三共用主解：master\_p23.py}
\lstinputlisting{../code/master_p23.py}

\subsection{问题二求解：problem2.py}
\lstinputlisting{../code/problem2.py}

\subsection{问题三求解：problem3.py}
\lstinputlisting{../code/problem3.py}

\subsection{问题四求解与效应隔离：problem4.py}
\lstinputlisting{../code/problem4.py}

\subsection{数值可信度检验：validate.py}
\lstinputlisting{../code/validate.py}

\subsection{报表写出：io\_out.py}
\lstinputlisting{../code/io_out.py}

\subsection{附件数据核对：data\_check.py}
\lstinputlisting{../code/data_check.py}

\subsection{模板结构核对：template\_check.py}
\lstinputlisting{../code/template_check.py}

\subsection{主程序入口：main.py}
\lstinputlisting{../code/main.py}
